\documentclass[11pt,a4paper]{altacv}

\geometry{left=1.5cm,right=1.5cm,top=1.cm,bottom=1.2cm,columnsep=1.5cm}

\usepackage[utf8]{inputenc}
\usepackage[T1]{fontenc}
\usepackage[french]{babel}
\usepackage{paracol}
\usepackage{xcolor}
\usepackage{hyperref}
\usepackage{fontawesome5}
\usepackage{setspace}

\definecolor{SlateGrey}{HTML}{2E2E2E}
\definecolor{LightGrey}{HTML}{666666}
\definecolor{DarkPastelRed}{HTML}{8AC0D9}
\definecolor{PastelRed}{HTML}{00B0DA}
\definecolor{GoldenEarth}{HTML}{B4AD0C}
\colorlet{name}{black}
\colorlet{tagline}{PastelRed}
\colorlet{heading}{DarkPastelRed}
\colorlet{headingrule}{GoldenEarth}
\colorlet{subheading}{PastelRed}
\colorlet{accent}{PastelRed}
\colorlet{emphasis}{SlateGrey}
\colorlet{body}{LightGrey}

\renewcommand{\familydefault}{\sfdefault}

\name{\Large Guillaume BOILEAU}
\tagline{Chef de projet Transport / Transition énergétique | Coordination, KPI, données \& pilotage}
\personalinfo{
  \email{guillaume.boileaupro@gmail.com}
  \phone{+33 6 59 94 85 84}
  \location{Alpes-Maritimes, France}
  \linkedin{guillaume-boileau-phd-891b81265}
}

\setlength{\parskip}{0.72\baselineskip}

\begin{document}
\makecvheader
\vspace{-0.65cm}

\cvsection{Lettre de motivation}
\vspace{-0.45cm}

{\color{accent}\textbf{Objet :}} \textit{Candidature -- Chef.fe de projet Transport PRODICEE -- Valbonne}

{\color{body}

Madame, Monsieur,

Je vous adresse ma candidature pour le poste de \textbf{Chef.fe de projet Transport PRODICEE}. L’ampleur du programme, son rôle d’évaluation technique et économique du dispositif des Certificats d'Économies d'Énergie, ainsi que sa contribution aux politiques publiques de transition énergétique correspondent à un type de mission dans lequel je souhaite m’investir : piloter, coordonner, analyser les données, produire des indicateurs fiables et faciliter la décision collective.

Mon parcours s’est construit à l’interface entre \textbf{R\&D, données, coordination transverse, validation, rédaction technique et systèmes complexes}. Docteur en physique, j’ai l’habitude de structurer des sujets techniques difficiles, de suivre des livrables, de produire des synthèses, de documenter des résultats et de coordonner des contributeurs autour d’objectifs communs. Ces compétences me semblent directement utiles pour piloter un volet du programme PRODICEE, animer des partenaires, suivre des KPI, consolider des contributions et restituer des résultats exploitables.

Mon expérience récente chez \textbf{VEV Platform Services France} constitue un lien concret avec la mobilité électrique. J’y ai travaillé sur une plateforme numérique CPMS connectée à des infrastructures de recharge de véhicules électriques. Cette expérience m’a confronté aux enjeux opérationnels des infrastructures de recharge : flux de données, incidents, fiabilité des échanges, comportements applicatifs, contraintes terrain et documentation technique. Elle me donne une compréhension pratique d’un environnement lié à la mobilité, aux services numériques associés et à l’exploitation de données opérationnelles.

Au \textbf{CNES}, j’ai travaillé sur le segment sol de la mission spatiale \textbf{LISA}, dans un environnement CNES/ESA associant chaînes de données mission, simulation, validation technique, coordination de contributeurs et revues collaboratives. J’y ai développé \textbf{CATpyLISA}, une plateforme Python destinée à structurer des workflows, comparer des modèles, valider des résultats et produire des éléments de synthèse. Cette expérience m’a donné une méthode solide pour cadrer un sujet, structurer les données, suivre les écarts, formaliser les résultats, préparer des revues et travailler avec des acteurs institutionnels et techniques.

À l’Universiteit Antwerpen, j’ai également développé des workflows de data science pour analyser des données capteurs, produire des indicateurs de performance, comparer des configurations et formuler des recommandations techniques. Cette expérience est directement transposable aux besoins d’un programme comme PRODICEE : croiser des données, évaluer des impacts, suivre des indicateurs, documenter les hypothèses et produire une analyse claire au service de la décision.

Je suis conscient que le poste demande une connaissance approfondie du dispositif CEE, du jeu d’acteurs transport et des partenaires nationaux comme le CEREMA ou l’ATEE. Mon parcours ne vient pas directement de ce secteur, mais j’apporte une capacité forte d’apprentissage, de structuration et de coordination. Je souhaite justement mettre mon expérience de pilotage technique, d’analyse de données et de restitution au service d’un programme public à impact, tout en montant rapidement en compétence sur le dispositif CEE Transport, les fiches standardisées, l’incitativité, les économies d’énergie réelles et les interactions avec les acteurs nationaux.

Pour le volet Transport de PRODICEE, je peux contribuer à plusieurs niveaux : suivi du planning et des actions, consolidation de KPI, analyse de données, coordination de partenaires, production de synthèses, préparation de comités, suivi de prestataires, documentation des résultats et communication avec les autres volets du programme. Mon expérience m’a appris à travailler avec rigueur, autonomie, diplomatie et esprit d’équipe dans des contextes où l’information doit être fiable, partagée et exploitable.

Rejoindre l’ADEME représente pour moi une opportunité de mettre mes compétences au service d’un projet directement utile à la transition écologique. Après des expériences en R\&D, spatial, data et mobilité électrique, je souhaite contribuer à des politiques publiques concrètes, en particulier sur des sujets mêlant transport, énergie, données, évaluation et coordination d’acteurs.

Je serais heureux de pouvoir échanger avec vous afin de vous présenter plus en détail mon parcours et la manière dont mon expérience en coordination, analyse de données, mobilité électrique, reporting et validation technique peut contribuer au volet Transport du programme PRODICEE.

Je vous prie d’agréer, Madame, Monsieur, l’expression de mes salutations distinguées.

\textbf{Guillaume Boileau}

}

\end{document}
